\section{Sintesis, Keterbatasan, dan Jawaban Rumusan Masalah}

Subbab ini menyintesis hasil keempat hipotesis lintas seluruh dataset, memosisikannya terhadap kebaruan yang diklaim, menjawab rumusan masalah, dan menyatakan keterbatasan.
Status ringkas H1-H4 disajikan pada Tabel \ref{tab:sintesis-hipotesis}.

\begin{longtable}{|>{\raggedright\arraybackslash}p{2cm}|>{\raggedright\arraybackslash}p{3cm}|>{\raggedright\arraybackslash}p{3cm}|>{\raggedright\arraybackslash}p{3cm}|}
\caption{Sintesis Status Hipotesis H1-H4 Lintas Dataset}
\label{tab:sintesis-hipotesis} \\
\hline
\textbf{Hipotesis} & \textbf{Dua dataset utama ($D=436$ \& $1.907$)} & \textbf{Tiga dataset kecil ($D=46$--$54$)} & \textbf{Smart Building IoT (eksternal)} \\
\hline
\endfirsthead

\multicolumn{4}{c}{{\bfseries \tablename\ \thetable{} -- lanjutan dari halaman sebelumnya}} \\
\hline
\textbf{Hipotesis} & \textbf{Dua dataset utama ($D=436$ \& $1.907$)} & \textbf{Tiga dataset kecil ($D=46$--$54$)} & \textbf{Smart Building IoT (eksternal)} \\
\hline
\endhead

\hline
\multicolumn{4}{r}{Lanjut ke halaman berikutnya} \\
\endfoot

\hline
\endlastfoot

% ==================== DATA BARIS ====================
H1 --- Kualitas & Unggul keluarga compact; setara iBOA; kalah moderat ACO/UBK pada D besar & Setara mayoritas (tumpul; duplikasi tinggi) & Melampaui sepele 5/5; setara iBOA \\
\hline
H2 --- Memori & Terkonfirmasi (datar $1{,}00\times$ vs $6{,}2$--$9{,}2\times$) & Konsisten (properti algoritma) & Konsisten \\
\hline
H3 --- Pembaruan & Terdukung (workforce stasioner); tak tampil di e-document (lantai derau) & Terbatas (kejenuhan) & Terdukung (mode sepadan; relearn 1/5) \\
\hline
H4 --- Warm-start & Terdukung penuh ($2{,}0$--$2{,}6\times$, kualitas naik) & Terdukung ($21$--$52\times$; kualitas campur) & Terdukung ($2{,}1$--$2{,}5\times$, kualitas naik) \\
\hline

\end{longtable}

Dua temuan utama penelitian adalah efisiensi memori dan pembaruan kebijakan. 
H2 menunjukkan bahwa memori CoDA-EDA tetap stabil terhadap $NP$, sedangkan BOA, ACO, dan UBK meningkat $6,2$–$9,2$ kali. 
H4 juga menunjukkan bahwa \textit{warm-start} mempercepat pembaruan kebijakan dengan kualitas yang tetap kompetitif. 
Sementara itu, H1 bergantung pada dimensi dan anggaran aturan, sedangkan H3 dipengaruhi tingkat \textit{drift} dan lantai derau JSD.

Hasil tersebut mendukung tiga unsur utama CoDA-EDA, yaitu pembaruan tanpa populasi eksplisit, pemodelan dependensi orde satu melalui Chow–Liu, dan pemisahan pembaruan parameter dari struktur. 
CoDA-EDA mampu mempertahankan pemodelan dependensi tanpa kehilangan sifat \textit{compact} serta mengurangi pembelajaran ulang struktur ketika distribusi stabil.

Dengan demikian, kontribusi utama CoDA-EDA terletak pada efisiensi memori dan pembaruan inkremental, bukan pada dominasi kualitas penambangan statis. 
Pengujian langsung pada Raspberry Pi 5 dan validasi lintas-\textit{generator} mendukung temuan tersebut, meskipun generalisasi masih terbatas oleh penggunaan satu dataset eksternal, data sintetis, serta beberapa hasil dengan jumlah replikasi terbatas.